% Generated from locale/tr/content/incompleteness/arithmetization-syntax/proofs-in-ax.tex; do not hand-edit.


\olfileid[tr]{inc}{art}{pax}
\olsection{Aksiyomatik \usetoken{P}{derivation}}

\begin{explain}
Aksiyomatik !!{derivation}s{}i aritmetikleştirmek için !!{derivation}s{}i
sayılarla temsil etmemiz gerekir. !!{derivation}s yalnızca !!{formula}
dizileri olduğundan en
açık yaklaşım, her !!{derivation}{}i içindeki !!{formula}s{}in kodlarından oluşan
dizinin koduyla kodlamaktır.
\end{explain}

\begin{defn}
$\delta$, $!A_1$, \dots,~$!A_n$ !!{formula}s{}inden oluşan aksiyomatik bir
!!{derivation} ise $\Gn{\delta}$
\[
\tuple{\Gn{!A_1}, \dots, \Gn{!A_n}}
\]
olur.
\end{defn}

\begin{ex}
  Şu çok basit !!{derivation}{}i ele alın:
  \begin{derivation}
  1. & $!B \lif (!B \lor !A)$ \\
  2. & $(!B \lif (!B \lor !A)) \lif (!A  \lif (!B \lif (!B \lor !A)))$\\
  3. & $!A  \lif (!B \lif (!B \lor !A))$
  \end{derivation}
  Bu !!{derivation}{}in Gödel numarası
  \begin{align*}
  \openTuple\,
    & \Gn{!B \lif (!B \lor !A)}, \\
    &\Gn{(!B \lif (!B \lor !A)) \lif (!A  \lif (!B \lif (!B \lor !A)))},\\
    & \Gn{!A  \lif (!B \lif (!B \lor !A))} \,\closeTuple
  \end{align*}
  olurdu.
\end{ex}

\begin{explain}
!!{derivation}s için bir temsil belirlediğimize göre bu !!{derivation}s{}i
ilkel özyinelemeli biçimde işleyebildiğimizi ve temel özellikleriyle
bağıntılarını da bu biçimde ifade edebildiğimizi göstermeliyiz. Bazı işlemler
basittir: Örneğin !!a{derivation}{}in $d$ Gödel numarası verildiğinde
$(d)_{\len{d}-1}$ bize son-!!{formula}{}ünün Gödel numarasını verir. Bazı
işlemler ise çok daha zordur. En azından bunların nasıl yapılacağını ana
hatlarıyla göstereceğiz. Amaç, ``$\delta$, $\Gamma$'dan $!A$'nın
!!a{derivation}{}idir'' bağıntısının $\delta$ ile $!A$'nın Gödel numaraları
üzerinde ilkel özyinelemeli olduğunu göstermektir.
\end{explain}

\begin{prop}
\ollabel{prop:followsby}
Aşağıdaki bağıntılar ilkel özyinelemelidir:
\begin{enumerate}
\item $!A$ bir aksiyomdur.
\item $\delta$'nın $i$'inci satırı modus ponens ile gerekçelendirilmiştir.
\item $\delta$'nın $i$'inci satırı \QR{} ile gerekçelendirilmiştir.
\item $\delta$ doğru bir !!{derivation}{}dir.
\end{enumerate}
\end{prop}

\begin{proof}
!!{formula}s{}in Gödel numaraları ile !!{derivation}s{}in Gödel numaraları
arasındaki karşılık gelen bağıntıların ilkel özyinelemeli olduğunu
göstermeliyiz.
\begin{enumerate}
\item Elimizde verilmiş bir aksiyom şemaları listesi vardır ve $!A$, ancak ve
  ancak bu şemalardan birinin verdiği biçimdeyse bir aksiyomdur. Şemaların
  listesi sonlu olduğundan, her aksiyom şeması için $!A$'nın o biçimde olup
  olmadığını ilkel özyinelemeli olarak sınayabildiğimizi göstermek yeterlidir.
  Örneğin
  \[
  !B \lif (!C \lif !B)
  \]
  aksiyom şemasını ele alın. `$($' ile $!B$'yi, `$\lif$' ile `$($'yi,
  $!C$'yi, `$\lif$' ile $!B$'yi ve son olarak `$))$'yi art arda
  birleştirdiğimizde $!A$'yı elde ettiğimiz !!{formula}s $!B$ ile $!C$ varsa,
  $!A$ bu aksiyom şemasının bir örneklemesidir. $!A$'nın Gödel numarası $n$
  üzerinde karşılık gelen özelliği sınayabiliriz; çünkü dizilerin
  birleştirilmesi ilkel özyinelemelidir. Ayrıca bağıntı geçerli olduğunda
  $!B$ ile $!C$, $!A$'nın alt-!!{formula}s{}i oldukları için Gödel numaraları
  $!A$'nın Gödel numarasından küçük olmalıdır. Dolayısıyla şöyle
  tanımlayabiliriz:
  \begin{multline*}
  \fn{IsAx}_{!B \lif (!C \lif !B)}(n) \defiff \bexists{b< n}{
    \bexists{c < n}{(\fn{Sent}(b) \land \fn{Sent}(c) \land {}}}\\ n =
  \Gn{(} \concat b \concat \Gn{\lif} \concat \Gn{(} \concat c \concat
  \Gn{\lif} \concat b \concat \Gn{))}
  \end{multline*}
  Her aksiyom şeması için böyle bir tanımımız varsa bunların ayrışımı,
  ``$n$ bir aksiyomun Gödel numarasıdır'' biçimindeki $\fn{IsAx}(n)$
  özelliğini tanımlar.
\item $\delta$'nın $i$'inci satırı, ancak ve ancak $j,k<i$ olan ve $j$'inci
  satırdaki !!{sentence} bir $!A$ formülü, $k$'ıncı satırdaki tümce
  $!A \lif !B$, $i$'inci satırdaki tümce de $!B$ olacak biçimde $j$ ile $k$
  satırları varsa modus ponens ile gerekçelendirilmiştir:
  \begin{multline*}
    \fn{MP}(d, i) \defiff \bexists{j < i}{\bexists{k < i}{}}\\
    (d)_k = \Gn{(} \concat (d)_j
      \concat \Gn{\lif} \concat (d)_i \concat \Gn{)}
  \end{multline*}
  Sınırlı niceleme, birleştirme ve $=$ ilkel özyinelemeli olduğundan bu, ilkel
  özyinelemeli bir bağıntı tanımlar.
\item $\delta$'nın bir satırı, $!B \lif \lforall[x][!A(x)]$ biçimindeyse,
  önceki bir satır bir $c$ !!{constant}{}i için $!B \lif !A(c)$ ise ve $c$,
  $!B$ içinde geçmiyorsa \QR{} ile gerekçelendirilmiştir. Bu, ancak ve ancak
  aşağıdakiler geçerliyse böyledir:
  \begin{enumerate}
  \item !!a{sentence}~$!B$ ve
  \item yalnızca tek bir $x$ değişkeninin serbest olduğu !!a{formula}~$!A(x)$
    vardır ve
  \item $i$'inci satır $!B \lif \lforall[x][!A(x)]$ içerir,
  \item bir $c$ sabiti için $j<i$ olan bir $j$ satırı
    $!B \lif \Subst{!A}{c}{x}$ içerir,
  \item ayrıca $c$, $!B$ içinde geçmez.
  \end{enumerate}
  Bunların tümü ilkel özyinelemeli biçimde sınanabilir. Çünkü $!B$, $!A(x)$
  ve $x$'in Gödel numaraları $i$'inci satırdaki formülün Gödel numarasından,
  $c$'nin Gödel numarası da $j$'inci satırdaki formülün Gödel numarasından
  küçüktür:
  \begin{multline*}
    \fn{QR}_1(d, i) \defiff \bexists{b < (d)_i}{\bexists{x <
        (d)_i}{\bexists{a < (d)_i}{\bexists{j<i}{\bexists{c < (d)_j}{(}}}}}
    \\ \fn{Var}(x) \land \fn{Const}(c) \land {} \\
    (d)_i = \Gn{(} \concat
    b \concat \Gn{\lif} \concat \Gn{\lforall} \concat x \concat a
    \concat \Gn{)} \land {}\\
    (d)_j = \Gn{(} \concat b \concat
    \Gn{\lif} \concat \fn{Subst}(a,c,x) \concat \Gn{)} \land {}\\ \fn{Sent}(b)
    \land \fn{Sent}(\fn{Subst}(a,c,x)) \land {} \bforall{k <
      \len{b}}{(b)_k \neq (c)_0})
  \end{multline*}
  Burada $x$ ile $c$'nin, ele alınan değişken ile sabitin terim olarak Gödel
  numaraları (yani simge kodları değil) olduğunu varsayıyoruz. $x$'in
  $!A(x)$'in tek serbest değişkeni olduğunu $\Subst{!A(x)}{c}{x}$'in
  !!a{sentence} olup olmadığını sınayarak belirleriz; $c$'nin $!B$ içinde
  geçmemesini ise $!B$'nin her simgesinin $c$'den farklı olmasını isteyerek
  sağlarız.

  \QR{} kuralının öteki sürümünü alıştırma olarak bırakıyoruz.
\item $d$, ancak ve ancak içindeki her satır bir aksiyomsa ya da modus ponens
  veya \QR{} ile gerekçelendirilmişse doğru bir !!{derivation}{}in Gödel
  numarasıdır. Dolayısıyla
  \[
  \fn{Deriv}(d) \defiff \bforall{i < \len{d}}{(\fn{IsAx}((d)_i) \lor
    \fn{MP}(d,i) \lor \fn{QR}(d, i))}
  \]
  olur.
\end{enumerate}
\end{proof}

\begin{prob}
Aşağıdaki bağıntıları \olref[inc][art][pax]{prop:followsby}'deki gibi
tanımlayın:
\begin{enumerate}
\item $\fn{IsAx}_{!A \lif (!B \lif (!A \land !B))}(n)$,
\item $\fn{IsAx}_{\lforall[x][!A(x)] \lif !A(t)}(n)$,
\item $\fn{QR}_{2}(d, i)$ (\QR{} kuralının öteki sürümü için).
\end{enumerate}
\end{prob}

\begin{prop}
$\Gamma$'nın ilkel özyinelemeli bir !!{sentence}s kümesi olduğunu varsayın.
O zaman ``$x$, $\Gamma$'dan $!A$'nın !!a{derivation}{}i olan $\delta$'nın kodudur
ve $y$, $!A$'nın Gödel numarasıdır'' ifadesini veren
$\Prf[\Gamma](x, y)$ bağıntısı ilkel özyinelemelidir.
\end{prop}

\begin{proof}
``$y \in \Gamma$'' ifadesinin ilkel özyinelemeli $R_\Gamma(y)$ yüklemiyle
verildiğini varsayın. $\Prf[\Gamma](x,y)$ bağıntısının ilkel özyinelemeli
olduğunu göstermeliyiz; burada $\Prf[\Gamma](x,y)$, ancak ve ancak $y$,
!!a{sentence}~$!A$'nın Gödel numarasıysa ve $x$, $\Gamma$'dan $!A$'nın
!!a{derivation}{}inin koduysa geçerlidir.

Önceki önermeye göre $x$'in doğru bir $\delta$ !!{derivation}{}inin kodu olması
durumunda geçerli olan $\fn{Deriv}(x)$ özelliği ilkel özyinelemelidir. Ancak
bu tanım, $\Gamma$ kümesini !!{derivation}{}deki satırları gerekçelendirmenin
ek bir yolu olarak hesaba katmamıştı. Bir satırın \QR{} ile gerekçelendirilip
gerekçelendirilmediğini sınayan ilkel özyinelemeli sınamamızda, $c$ sabitinin
$\Gamma$'da geçmesine izin verilmediği koşulunu da hesaba katmamıştık.
Tanımımızı $\Gamma$'yı doğrudan hesaba katacak biçimde değiştirmek mümkündür;
ancak $\fn{Deriv}$'i ve tümdengelim teoremini kullanmak daha kolaydır.
$\Gamma \Proves !A$ olması, ancak ve ancak $!B_1$, \dots, $!B_n \in
\Gamma$ olan sonlu bir !!{sentence}s listesi için
$\{!B_1, \dots, !B_n\} \Proves !A$ olması demektir. Tümdengelim teoremine
göre bu, $\Proves (!B_1 \lif (!B_2 \lif
\cdots (!B_n \lif !A)\cdots))$ olduğunda geçerlidir. Gödel numarası $z$ olan
!!a{sentence}{}nin bu biçimde olup olmadığı ilkel özyinelemeli biçimde
sınanabilir. Böylece $x$'i, Gödel numarası $y$ olan !!{sentence}{}nin
\emph{$\Gamma$'dan} !!a{derivation}{}inin Gödel numarası olarak görmek yerine,
yukarıdaki biçimde iç içe geçmiş bir koşullunun $\emptyset$'tan
!!a{derivation}{}inin Gödel numarası olarak görürüz.

Önce bir !!{sentence}s dizimiz varsa bütün bu tümceleri önbileşen, verilmiş
!!{sentence}{}yi da artbileşen yapan koşulluyu ilkel özyinelemeli biçimde
kurabiliriz:
\begin{align*}
  \fn{hCond}(s, y, 0) & = y \\
  \fn{hCond}(s, y, n+1) & = \Gn{(} \concat (s)_{n} \concat \Gn{\lif}
  \concat \fn{hCond}(s, y, n) \concat \Gn{)}\\
  \fn{Cond}(s, y) & = \fn{hCond}(s, y, \len{s})\\
  \intertext{Böylece $\Prf[\Gamma](x, y)$'yi şöyle tanımlayabiliriz:}
  \Prf[\Gamma](x, y) & \defiff \bexists{s < \fn{sequenceBound}(x,x)}{(} \\
  &\qquad (x)_{\len{x}-1} = \fn{Cond}(s,y) \land {} \\
  &\qquad\bforall{i<\len{s}}{(s)_i \in \Gamma} \land {} \\
  &\qquad\fn{Deriv}(x))
\end{align*}
$s$ üzerindeki sınır şöyle elde edilir: Her $(s)_i$, !!{derivation}{}in son
satırındaki bir alt-!!{formula}{}ün Gödel numarasıdır; dolayısıyla
$(x)_{\len{x}-1}$'den küçüktür. $!B \in \Gamma$ biçimindeki önbileşenlerin
sayısı olan $s$'nin uzunluğu da $x$'in son satırının uzunluğundan
küçüktür.
\end{proof}

